\documentclass[12pt]{article}
\usepackage[utf8]{inputenc}
\usepackage[T1]{fontenc}
\usepackage[margin=1in]{geometry}
\usepackage{setspace}
\usepackage{xcolor}
\usepackage[colorlinks=true, urlcolor=blue!60!black, linkcolor=blue!60!black]{hyperref}
\setlength{\parindent}{2em}
\setlength{\parskip}{1.25em}
\renewcommand{\baselinestretch}{1.0}
\begin{document}
\begin{center}
{\LARGE \textit{District of Columbia v. Heller}, 554 U.S. 570 (2008)\par}
\vspace{0.8em}
{\large Case Brief\par}
\end{center}
\vspace{1em}

\section{Memory Jogger}
In a landmark 5--4 originalist decision, the Supreme Court held for the first time that the Second Amendment protects an individual right to possess and carry firearms unconnected with militia service, striking down Washington, D.C.'s handgun ban and trigger-lock mandate as unconstitutional infringements on self-defense in the home.

\section{Detailed Case Facts}
In 1975, responding to high rates of urban gun crime, the District of Columbia enacted the Firearms Control Regulations Act of 1975, creating one of the strictest municipal firearms regimes in the United States. Three interlocking provisions of the D.C. Code were at issue: (1) D.C. Code \S\S\ 7-2501.01(12), 7-2502.01(a), and 7-2502.02(a)(4), which made it a crime to carry an unregistered firearm and prohibited the registration of handguns (effectively banning handgun possession by civilians); (2) D.C. Code \S\S\ 22-4504(a) and 22-4506, which barred carrying an unlicensed handgun (even inside one's own home) while authorizing the police chief to issue one-year licenses; and (3) D.C. Code \S\ 7-2507.02, which required residents to keep lawfully owned firearms---such as registered rifles and shotguns---``unloaded and disassembled or bound by a trigger lock or similar device'' unless located at a place of business or used for lawful recreational activities, with no statutory exception for self-defense (\href{https://www.law.cornell.edu/supct/html/07-290.ZO.html}{Cornell LII}; \href{https://www.courtlistener.com/opinion/145777/district-of-columbia-v-heller/}{CourtListener}). Violations carried criminal penalties of up to one year in prison for a first offense and five years for a second offense (D.C. Code \S\ 7-2507.06).

Respondent Dick Anthony Heller was a District of Columbia special police officer authorized to carry a handgun while on duty at the Federal Judicial Center. Heller applied for a registration certificate to keep a handgun at his home in a high-crime Washington neighborhood for personal protection, but the District denied his application under its statutory ban. Because the District of Columbia is a federal enclave rather than a state, the case directly presented the meaning of the Second Amendment---``A well regulated Militia, being necessary to the security of a free State, the right of the people to keep and bear Arms, shall not be infringed''---without requiring the Court to decide whether the Amendment was incorporated against the states through the Fourteenth Amendment (a question resolved two years later in \href{https://supreme.justia.com/cases/federal/us/561/742/}{\textit{McDonald v. City of Chicago}}, 561 U.S. 742 (2010)).

Historically and politically, \textit{Heller} arrived after nearly seventy years of lower-court consensus reading \textit{United States v. Miller}, 307 U.S. 174 (1939), as tying the Second Amendment to militia service, disrupted only recently by the Fifth Circuit's individual-rights decision in \textit{United States v. Emerson}, 270 F.3d 203 (5th Cir. 2001), and decades of concerted legal scholarship and conservative constitutional advocacy.

\section{Procedural History}
Six D.C. residents, led initially by Shelly Parker and including Dick Heller, filed suit in the United States District Court for the District of Columbia seeking to enjoin enforcement of the handgun registration bar, the in-home licensing restriction, and the trigger-lock/disassembly mandate on Second Amendment grounds. In 2004, the District Court dismissed the complaint, holding that the Second Amendment did not confer an individual right to possess firearms outside militia service (\textit{Parker v. District of Columbia}, 311 F. Supp. 2d 103 (D.D.C. 2004)) (\href{https://www.law.cornell.edu/supct/html/07-290.ZO.html}{Cornell LII}).

On appeal, the United States Court of Appeals for the District of Columbia Circuit held that five of the six plaintiffs lacked Article III standing because they had not suffered a concrete injury, but that Heller had standing because his handgun registration application had been formally denied. Reaching the merits in a 2--1 decision, the D.C. Circuit reversed, holding that the Second Amendment protects an individual right to possess firearms unconnected with militia service and that D.C.'s total handgun ban and functional-firearm prohibition violated that right (\textit{Parker v. District of Columbia}, 478 F.3d 370 (D.C. Cir. 2007)). The D.C. Circuit directed entry of summary judgment for Heller. The District of Columbia petitioned for certiorari, which the Supreme Court granted on November 20, 2007 (\href{https://www.oyez.org/cases/2007/07-290}{Oyez}). Following oral argument on March 18, 2008, the Supreme Court affirmed the D.C. Circuit's judgment on June 26, 2008 (\href{https://supreme.justia.com/cases/federal/us/554/570/}{Justia}).

\section{Judicial Votes}
\begin{itemize}
  \item \textbf{Majority (5--4):} Justice Scalia delivered the opinion of the Court, joined by Chief Justice Roberts and Justices Kennedy, Thomas, and Alito.
  \item \textbf{Concurrences:} None.
  \item \textbf{Dissents (4):}
    \begin{itemize}
      \item Justice Stevens filed a dissenting opinion, joined by Justices Souter, Ginsburg, and Breyer.
      \item Justice Breyer filed a separate dissenting opinion, joined by Justices Stevens, Souter, and Ginsburg.
    \end{itemize}
\end{itemize}

\section{Holding}
\textbf{Judgment:} Affirmed (5--4). The District of Columbia's total ban on handgun possession in the home and its requirement that lawful firearms in the home be kept unloaded and disassembled or bound by a trigger lock violate the Second Amendment. Assuming Heller is not otherwise disqualified, the District must allow him to register his handgun and issue him a license to carry it in his home.

\textbf{Rule Announced:} The Second Amendment guarantees a pre-existing individual right to keep and bear arms unconnected with service in an organized militia, and to use arms for traditionally lawful purposes, with self-defense in the home at its core. Under any standard of constitutional scrutiny applicable to enumerated rights (excluding rational-basis review and freestanding interest-balancing), a categorical ban on an entire class of arms in common use for lawful self-defense---handguns---and a mandate rendering firearms inoperable in the home are unconstitutional (\href{https://www.law.cornell.edu/supct/html/07-290.ZS.html}{Cornell LII Syllabus}).

\textbf{Holding vs. Dicta:} The strict holding invalidates D.C.'s categorical in-home handgun ban and operable-firearm prohibition (without deciding the validity of D.C.'s licensing regime, which Heller conceded at oral argument was permissible if administered non-arbitrarily, and without deciding Fourteenth Amendment incorporation). Crucially, Part III's famous ``safe harbor'' passage---stating that nothing in the opinion casts doubt on ``longstanding prohibitions on the possession of firearms by felons and the mentally ill,'' bans on carrying in ``sensitive places such as schools and government buildings,'' conditions on commercial sales, or bans on ``dangerous and unusual weapons'' such as M-16 rifles---is widely regarded as influential dicta, though lower courts and subsequent Supreme Court decisions have treated it as highly authoritative (\href{https://www.heritage.org/courts/report/long-standing-and-presumptively-lawful-hellers-dicta-vs-history-and-dicta}{Heritage Foundation}).

\section{Analysis of Opinions}

\subsection{Majority Opinion (Justice Scalia)}
Justice Scalia grounded the Court's opinion in public-meaning originalism, asking how ordinary citizens of the founding generation understood the Second Amendment's words (\href{https://www.law.cornell.edu/supct/html/07-290.ZO.html}{Cornell LII}). Structurally, Scalia divided the Amendment into a ``prefatory clause'' (``A well regulated Militia, being necessary to the security of a free State'') and an ``operative clause'' (``the right of the people to keep and bear Arms, shall not be infringed''). Invoking 18th- and 19th-century statutory construction treatises, he held that a prefatory clause announces a purpose and may resolve ambiguity, but does not limit or expand an unambiguous operative command.

Analyzing the operative clause first, Scalia reasoned that ``the right of the people'' consistently refers to individual rights belonging to all members of the political community across the First, Fourth, and Ninth Amendments (citing \textit{United States v. Verdugo-Urquidez}, 494 U.S. 259 (1990)), in contrast to ``the militia,'' which was a subset of able-bodied men. Drawing on 18th-century dictionaries (Johnson, Cunningham, Webster), Blackstone, and early state constitutions, Scalia concluded that ``Arms'' meant bearable weapons generally (including modern arms not existing in 1791), ``keep arms'' simply meant possessing weapons, and ``bear arms'' meant carrying weapons for offensive or defensive confrontation---bearing an exclusively military idiom only when followed by ``against.''

Turning to the prefatory clause, Scalia argued that ``the Militia'' meant the unorganized pool of all able-bodied male citizens, and ``well-regulated'' meant properly disciplined and trained. The two clauses fit ``perfectly'' because tyrants (such as the Stuart kings before the 1689 English Bill of Rights) destroyed citizen militias by disarming the general populace. Thus, preserving the citizens' militia was the reason the pre-existing individual right of self-defense was codified, even though self-defense remained the right's ``central component.'' Scalia corroborated this reading with analogous pre- and post-1789 state constitutional provisions, early commentators (St. George Tucker, William Rawle, Joseph Story), antebellum state cases (\textit{Nunn v. State}, \textit{State v. Chandler}), and Reconstruction-era legislation (the Freedmen's Bureau Act of 1866) and commentary (Thomas Cooley).

Addressing precedent, Scalia contended that neither \textit{United States v. Cruikshank}, 92 U.S. 542 (1876), nor \textit{Presser v. Illinois}, 116 U.S. 252 (1886), foreclosed an individual right. He narrowly cabined \textit{United States v. Miller}, 307 U.S. 174 (1939)---noting the respondent had filed no brief and did not appear---as holding only that the Second Amendment does not protect weapons ``not typically possessed by law-abiding citizens for lawful purposes,'' such as short-barreled shotguns, while protecting arms ``in common use at the time.'' Finally, in Parts III and IV, Scalia affirmed that the right is ``not unlimited,'' listed ``presumptively lawful'' regulations, rejected rational-basis review and Justice Breyer's ``judge-empowering `interest-balancing inquiry,'\,'' and struck down D.C.'s handgun ban and trigger-lock requirement without selecting a formal tier of scrutiny.

\subsection{Dissenting Opinion (Justice Stevens)}
Justice Stevens, joined by Justices Souter, Ginsburg, and Breyer, argued that the Second Amendment was adopted solely to protect the right of the people of the several States to maintain a well-regulated militia against federal disarmament and standing armies---not to constitutionalize private self-defense or curtail legislative regulation of civilian gun ownership (\href{https://www.law.cornell.edu/supct/html/07-290.ZD.html}{Cornell LII}).

Stevens attacked the majority's textual, historical, and precedential analysis on four fronts. First, he faulted Scalia for inverting ordinary interpretation by relegating the preamble to an afterthought, violating \textit{Marbury v. Madison}'s rule against surplusage. Second, relying on linguistics amici and founding-era usage, Stevens maintained that ``keep and bear Arms'' was a unitary military phrase---quoting state militia statutes requiring men to ``keep'' arms ready for muster and noting that unadorned ``bear arms'' idiomatically meant military service. Third, Stevens emphasized the Amendment's drafting history: James Madison modeled the text on Virginia's military-focused proposal and included a conscientious-objector clause (later deleted in the Senate), while rejecting alternative state proposals from Pennsylvania and New Hampshire that expressly protected hunting or personal self-defense. Fourth, Stevens criticized the majority's reliance on the 1689 English Bill of Rights (which bound only the Crown, not Parliament), post-enactment 19th-century commentary, and post-Civil War history, and accused the Court of abandoning \textit{stare decisis} by overturning seven decades of settled reliance on \textit{Miller} and \textit{Lewis v. United States}, 445 U.S. 55 (1980).

\subsection{Dissenting Opinion (Justice Breyer)}
Justice Breyer, joined by Justices Stevens, Souter, and Ginsburg, agreed with Stevens's militia-centered reading but argued independently that even assuming the Second Amendment protects an individual self-defense interest, D.C.'s handgun ban was constitutional (\href{https://www.law.cornell.edu/constitution-conan/amendment-2/the-heller-decision-and-individual-right-to-firearms}{Constitution Annotated}). First, Breyer surveyed colonial urban fire-safety and firearm-discharge regulations in Boston, Philadelphia, and New York (such as a 1783 Boston law barring loaded firearms inside buildings) to argue that founding-era legislatures imposed proportionate burdens on urban gun possession. Second, Breyer rejected strict scrutiny as unworkable because virtually every gun regulation serves a compelling interest in saving lives, and instead proposed a pragmatic ``interest-balancing inquiry'' asking ``whether the statute burdens a protected interest in a way or to an extent that is out of proportion to the statute's salutary effects upon other important governmental interests'' (\href{https://www.judiciary.senate.gov/download/heller-v-district-of-columbia_-appendix-and-conclusion}{D.C. Cir. citing 554 U.S. at 689--690}). Emphasizing empirical evidence of urban handgun homicides and legislative competence in resolving contested social-science data, Breyer concluded that D.C.'s tailored urban restriction---which still permitted rifles and shotguns---was a proportionate, constitutionally valid response to gun violence.

\section{Examples: Future Applications}
\begin{enumerate}
  \item \textbf{Same-Side Application (State and Municipal Handgun Bans):} A major city such as Chicago enacts an ordinance mirroring D.C.'s registration bar, effectively prohibiting law-abiding residents from keeping handguns in their homes for self-defense. Once the Second Amendment is incorporated against the states via the Fourteenth Amendment, \textit{Heller} compels invalidation of the ban, as the Supreme Court held in \href{https://supreme.justia.com/cases/federal/us/561/742/}{\textit{McDonald v. City of Chicago}}, 561 U.S. 742 (2010).
  \item \textbf{Same-Side Application (Discretionary ``Proper Cause'' Public-Carry Licensing):} A state prohibits law-abiding citizens from carrying a handgun in public for self-defense unless they prove a ``special need'' beyond that of the general public. Extending \textit{Heller}'s textual-and-historical methodology and rejection of interest-balancing beyond the home, such discretionary licensing violates the Second Amendment, as the Court held in \href{https://supreme.justia.com/cases/federal/us/597/20-843/}{\textit{New York State Rifle \& Pistol Ass'n v. Bruen}}, 597 U.S. 1 (2022).
  \item \textbf{Opposite-Side Application (Temporary Disarmament of Domestic-Violence Offenders):} Federal law (18 U.S.C. \S\ 922(g)(8)) bars an individual subject to a court-issued domestic-violence restraining order---entered after a judicial finding that he poses a credible threat to an intimate partner's physical safety---from possessing firearms. Consistent with \textit{Heller}'s recognition that the right is ``not unlimited'' and historical traditions of disarming dangerous individuals, the prohibition is constitutional, as the Court held in \href{https://www.supremecourt.gov/opinions/23pdf/22-915_8o6b.pdf}{\textit{United States v. Rahimi}}, 602 U.S. 680 (2024).
  \item \textbf{Opposite-Side Application (Ban on Civilian Possession of Fully Automatic Machine Guns or Short-Barreled Shotguns):} Congress prohibits private possession of newly manufactured M-16 rifles or unregistered sawed-off shotguns. Under \textit{Heller}'s reading of \textit{Miller} and the historical tradition barring ``dangerous and unusual weapons'' not in common use for lawful self-defense, such restrictions remain valid.
  \item \textbf{Fence-Sitter Application (State Bans on Semiautomatic ``Assault Weapons'' and Large-Capacity Magazines):} A state bans civilian sale and possession of AR-15-style semiautomatic rifles and magazines holding more than ten rounds. Challengers argue under \textit{Heller} that millions of law-abiding Americans own such rifles and magazines, making them ``in common use'' for lawful purposes and thus categorically protected; defenders argue that their military-grade lethality renders them ``dangerous and unusual'' weapons akin to M-16s and that bans leave ample self-defense alternatives (including handguns) untouched.
\end{enumerate}

\section{Critique}
\textit{Heller} is simultaneously hailed as the ``triumph of originalism'' and criticized from both left and right as a cautionary tale about judicial activism and the limits of historical jurisprudence.

\textbf{Logical Weaknesses and Internal Contradictions:} Critics highlight several internal tensions in Justice Scalia's opinion. First, having anchored the right in the textual argument that ``the people'' encompasses all members of the political community just as in the First and Fourth Amendments, Scalia narrows the right at the end of the opinion to ``law-abiding, responsible citizens'' and summarily approves felon-in-possession bans---even though felons retain First and Fourth Amendment rights. Second, \textit{Heller}'s ``common use'' test creates an awkward circularity: whether a weapon is protected depends on how widely civilians possess it, which in turn depends on whether legislatures have historically banned it. Third, the majority's laundry list of ``presumptively lawful'' regulations in Part III (bans for felons and the mentally ill, sensitive-place bans, commercial-sale conditions) is asserted by ipse dixit without any founding-era historical proof, undercutting the opinion's methodological claim to strict originalism (\href{https://repository.uclawsf.edu/cgi/viewcontent.cgi?article=1426&context=faculty_scholarship}{UC Law SF}; \href{https://www.heritage.org/courts/report/long-standing-and-presumptively-lawful-hellers-dicta-vs-history-and-dicta}{Heritage Foundation}).

\textbf{Progressive and Pragmatic Critique:} Progressive scholars and historians---alongside prominent conservative pragmatists---argue that \textit{Heller} engaged in selective ``law-office history'' and undermined democratic self-governance. Yale Law Professor Reva Siegel argues that \textit{Heller}'s originalism actually reflects late-20th-century ``popular constitutionalism'' forged by the modern gun-rights movement rather than 1791 public meaning (\href{https://law.yale.edu/sites/default/files/documents/pdf/Faculty/Siegel_DeadOrAliveOriginalismAsPopularConstitutionalismInHeller.pdf}{Siegel, Yale Law}). Judges Richard Posner and J. Harvie Wilkinson III---both conservative appellate jurists---famously attacked \textit{Heller} as the conservative equivalent of \textit{Roe v. Wade}: Posner faulted Scalia's ``faux originalism'' for ignoring the Amendment's militia preamble and disabling local experimentation against urban gun violence, while Wilkinson argued that when history is genuinely ambiguous, judicial restraint demands deferring to elected legislatures rather than constitutionalizing a divisive social policy (\href{https://scholarship.law.nd.edu/cgi/viewcontent.cgi?article=4582&context=ndlr}{Notre Dame Law Review}).

\textbf{Originalist and Textualist Perspectives:} Many originalists celebrate \textit{Heller} as a watershed victory for original public meaning over living constitutionalism and purposive balancing, praising Scalia's refusal to let a prefatory clause erase an operative right or to subordinate an enumerated guarantee to Justice Breyer's ad hoc proportionality test. Yet even sympathetic originalist scholars (such as Nelson Lund) have criticized \textit{Heller} for two opposing sins: failing to articulate a clear doctrinal standard of review (which invited lower courts for fourteen years to apply watered-down intermediate scrutiny until \textit{Bruen} repudiated it) and inserting unsubstantiated, non-originalist ``presumptively lawful'' exceptions in Part III that diluted the Amendment's text to secure Justice Kennedy's decisive fifth vote.

\section{Key Quotations}
\begin{quote}
``Putting all of these textual elements together, we find that they guarantee the individual right to possess and carry weapons in case of confrontation.''
\end{quote}
\noindent --- Justice Scalia, Opinion of the Court, 554 U.S. at 592 (\href{https://www.law.cornell.edu/supct/html/07-290.ZO.html}{Cornell LII}).

\begin{quote}
``Although we do not undertake an exhaustive historical analysis today of the full scope of the Second Amendment, nothing in our opinion should be taken to cast doubt on longstanding prohibitions on the possession of firearms by felons and the mentally ill, or laws forbidding the carrying of firearms in sensitive places such as schools and government buildings, or laws imposing conditions and qualifications on the commercial sale of arms.''
\end{quote}
\noindent --- Justice Scalia, Opinion of the Court, 554 U.S. at 626--627 (\href{https://www.law.cornell.edu/supct/html/07-290.ZO.html}{Cornell LII}).

\begin{quote}
``We know of no other enumerated constitutional right whose core protection has been subjected to a freestanding `interest-balancing' approach. The very enumeration of the right takes out of the hands of government---even the Third Branch of Government---the power to decide on a case-by-case basis whether the right is really worth insisting upon. A constitutional guarantee subject to future judges' assessments of its usefulness is no constitutional guarantee at all.''
\end{quote}
\noindent --- Justice Scalia, Opinion of the Court, 554 U.S. at 634 (\href{https://www.law.cornell.edu/supct/html/07-290.ZO.html}{Cornell LII}).

\begin{quote}
``But the enshrinement of constitutional rights necessarily takes certain policy choices off the table. These include the absolute prohibition of handguns held and used for self-defense in the home.''
\end{quote}
\noindent --- Justice Scalia, Opinion of the Court, 554 U.S. at 636 (\href{https://www.law.cornell.edu/supct/html/07-290.ZO.html}{Cornell LII}).

\begin{quote}
``The Second Amendment was adopted to protect the right of the people of each of the several States to maintain a well-regulated militia. It was a response to concerns raised during the ratification of the Constitution that the power of Congress to disarm the state militias and create a national standing army posed an intolerable threat to the sovereignty of the several States.''
\end{quote}
\noindent --- Justice Stevens, dissenting, 554 U.S. at 637 (\href{https://www.law.cornell.edu/supct/html/07-290.ZD.html}{Cornell LII}).


\end{document}
